\subsection{PAC-Bayes and Information-Theoretic Bounds}

PAC-Bayes bounds analyze randomized predictors by comparing a data-dependent posterior distribution on hypotheses to a fixed prior.

\begin{definition} \label{def:sl-randomized-predictor} \hlt{Randomized Predictor}\\
Let $\mathcal{H}$ be a hypothesis space and $Q\in\mathcal{P}(\mathcal{H})$. The \hlt{randomized predictor} $Q$ samples $h\sim Q$ before prediction, with risk $R(Q)\defeq\mathbb{E}_{h\sim Q}R(h)$ and empirical risk $\widehat{R}_S(Q)\defeq\mathbb{E}_{h\sim Q}\widehat{R}_S(h)$.
\end{definition}

\begin{theorem} \label{thm:sl-pac-bayes-kl-form} \hlt{PAC-Bayes Bound Form}\\
Let $S_n=(Z_1,\ldots,Z_n)$ be i.i.d., let $\ell(h,Z)\in[0,1]$, and let $\Pi\in\mathcal{P}(\mathcal{H})$ be independent of $S_n$.\\
For $Q\ll\Pi$, set $\operatorname{KL}(Q\Vert\Pi)\defeq\mathbb{E}_Q\log\frac{dQ}{d\Pi}$ and $\operatorname{kl}(a,b)\defeq a\log\frac{a}{b}+(1-a)\log\frac{1-a}{1-b}$.
\begin{enumerate}[label=(\roman*)]
\setlength{\itemsep}{0pt}
\item With probability at least $1-\delta$, every posterior $Q\ll\Pi$ satisfies $\operatorname{kl}(\widehat{R}_{S_n}(Q),R(Q))\leq\frac{\operatorname{KL}(Q\Vert\Pi)+\log\left(\frac{n+1}{\delta}\right)}{n}$.
\item On the same event, $R(Q)\leq\widehat{R}_{S_n}(Q)+\sqrt{\frac{\operatorname{KL}(Q\Vert\Pi)+\log\left(\frac{n+1}{\delta}\right)}{2n}}$ for every $Q\ll\Pi$.
\end{enumerate}
\begin{proof}
\begin{enumerate}[label=(\roman*)]
\setlength{\itemsep}{0pt}
\item For fixed $h$, the bounded-loss kl moment bound gives $\mathbb{E}_{S_n}\exp(n\operatorname{kl}(\widehat{R}_{S_n}(h),R(h)))\leq n+1$.\\
Thus $\mathbb{E}_{S_n}\mathbb{E}_{h\sim\Pi}\exp(n\operatorname{kl}(\widehat{R}_{S_n}(h),R(h)))\leq n+1$, so Markov gives,\\
with probability at least $1-\delta$, $\mathbb{E}_{\Pi}\exp(n\operatorname{kl}(\widehat{R}_{S_n}(h),R(h)))\leq\frac{n+1}{\delta}$.\\
On this event, the Donsker-Varadhan inequality with $\varphi(h)\defeq n\operatorname{kl}(\widehat{R}_{S_n}(h),R(h))$ gives\\
$\mathbb{E}_Q\varphi\leq\operatorname{KL}(Q\Vert\Pi)+\log\mathbb{E}_{\Pi}\exp\varphi$  $\forall Q\ll\Pi$. By convexity of $\operatorname{kl}$, $n\operatorname{kl}(\widehat{R}_{S_n}(Q),R(Q))\leq\mathbb{E}_Q\varphi$, proving (i).
\item If $R(Q)\geq\widehat{R}_{S_n}(Q)$, Pinsker gives $\operatorname{kl}(\widehat{R}_{S_n}(Q),R(Q))\geq2(R(Q)-\widehat{R}_{S_n}(Q))^2$.\\
if $R(Q)<\widehat{R}_{S_n}(Q)$, the claim is immediate. Combining with (i) gives (ii).
\end{enumerate}
\hlt{Structural Commentary}: PAC-Bayes replaces class size by posterior-prior divergence.\\
\hlt{Key Insight}: A prior-average exponential tail plus $Q\ll\Pi$ yields all posteriors at cost $\operatorname{KL}(Q\Vert\Pi)$.
\end{proof}
\end{theorem}

\begin{remark} \label{rem:sl-pac-bayes-prior-posterior} \hlt{Prior and Posterior}\\
In PAC-Bayes, $\Pi$ need not be subjective; it anchors the complexity term $\operatorname{KL}(Q\Vert\Pi)$.\\
Smaller $\operatorname{KL}(Q\Vert\Pi)$ means the data-dependent posterior $Q$ stays closer to the fixed prior.
\end{remark}

\begin{definition} \label{def:sl-information-generalization-gap} \hlt{Information-Theoretic Generalization}\\
Let $S_n=(Z_1,\ldots,Z_n)$, $W=A(S_n)$ be learned output. \hlt{Information-theoretic generalization} bounds control $\operatorname{gen}(A)\defeq\mathbb{E}[R(W)-\widehat{R}_{S_n}(W)]$ by dependence measures such as $I(W;S_n)$.
\end{definition}
